\documentclass[a4paper,12pt]{article}
\usepackage[T2A]{fontenc}
\usepackage[utf8]{inputenc}
\usepackage[russian]{babel}
\usepackage{amsmath,amssymb,mathtools}
\usepackage{helvet}
\renewcommand{\familydefault}{\sfdefault}
\usepackage{icomma}
\usepackage{geometry}
\geometry{left=24mm,right=24mm,top=23mm,bottom=24mm}
\usepackage{microtype}
\usepackage{xcolor}
\usepackage{tikz}
\usetikzlibrary{calc,arrows.meta,positioning,shapes.geometric}
\usepackage{booktabs,tabularx}
\usepackage{enumitem}
\usepackage{hyperref}
\usepackage{fancyhdr}
\usepackage{setspace}
\usepackage{needspace}
\usepackage{array}

\definecolor{accent}{HTML}{147A7C}
\definecolor{darkaccent}{HTML}{155457}
\definecolor{textgray}{HTML}{30383A}
\definecolor{muted}{HTML}{687576}
\definecolor{soft}{HTML}{D5E7E5}
\hypersetup{colorlinks=true,linkcolor=darkaccent,urlcolor=darkaccent,pdfauthor={Артём Александрович Лёвин},pdftitle={Чек-лист. НОД и НОК: олимпиадные приёмы. 12.10.26}}
\setstretch{1.28}
\setlength{\parindent}{0pt}
\setlength{\parskip}{5pt}
\setlist[itemize]{leftmargin=6mm,itemsep=2pt,topsep=3pt}
\setlist[enumerate]{leftmargin=8mm,itemsep=5pt,topsep=3pt}
\setlength{\headheight}{15pt}
\pagestyle{fancy}
\fancyhf{}
\fancyhead[L]{\small\color{muted}Математика \,$\cdot$\, олимпиадная подготовка}
\fancyhead[R]{\small\color{muted}12.10.2026}
\fancyfoot[C]{\small\color{muted}\thepage}
\renewcommand{\headrulewidth}{0pt}
\newcommand{\topic}[1]{\needspace{6\baselineskip}\vspace{9pt}{\Large\bfseries\color{darkaccent}#1}\par\vspace{2pt}\noindent\color{soft}\rule{\linewidth}{0.6pt}\color{textgray}\par\vspace{5pt}}
\newcommand{\subtopic}[1]{\needspace{3\baselineskip}\vspace{4pt}{\large\bfseries\color{darkaccent}#1}\par}
\newcommand{\idea}[1]{\noindent\makebox[6mm][l]{\color{accent}$\square$}\begin{minipage}[t]{\dimexpr\linewidth-7mm}\raggedright #1\end{minipage}\par\vspace{3pt}}
\newcommand{\formula}[1]{\[\boxed{#1}\]}
\newcommand{\gcdop}{\mathop{\text{НОД}}\nolimits}
\newcommand{\lcmop}{\mathop{\text{НОК}}\nolimits}
\newcommand{\pair}[2]{(#1;\,#2)}
\AtBeginDocument{\color{textgray}}

\begin{document}
\thispagestyle{empty}
\begin{center}
\vspace*{28mm}
{\small\color{accent}\bfseries ОЛИМПИАДНАЯ МАТЕМАТИКА}\par
\vspace{13mm}
{\fontsize{31}{37}\selectfont\bfseries\color{darkaccent}Чек-лист}\par
\vspace{8mm}
{\fontsize{21}{28}\selectfont\bfseries НОД и НОК:}\par
\vspace{2mm}
{\fontsize{19}{26}\selectfont олимпиадные приёмы}\par
\vspace{13mm}
{\large Разложение на простые множители}\par
{\large Алгоритм Евклида и взаимная простота}\par
{\large Задачи с параметром и восстановление чисел}\par
\vspace{17mm}
{\color{accent}\rule{38mm}{0.7pt}}\par
\vspace{9mm}
{\large\bfseries Ксения Клыкова}\par
\vspace{3mm}
{\normalsize 12 октября 2026 года}\par
\vfill
{\small\color{muted}Лёвин Артём Александрович}\par
{\small\color{darkaccent}\bfseries эксклюзивно для Ксении Клыковой}\par
\vspace{20mm}
\end{center}
\begin{tikzpicture}[remember picture,overlay]
  \draw[accent!55,line width=1.1pt] ([xshift=25mm,yshift=25mm]current page.south west)
    .. controls ([xshift=65mm,yshift=41mm]current page.south west)
    and ([xshift=-65mm,yshift=11mm]current page.south east)
    .. ([xshift=-24mm,yshift=27mm]current page.south east);
\end{tikzpicture}
\clearpage

\topic{1. Простые множители: минимум и максимум}
\idea{Каждое натуральное число больше $1$ можно представить как произведение степеней простых чисел. Для сравнения разложений удобно использовать \emph{один и тот же набор} простых оснований; отсутствующий множитель имеет степень $0$.}
\idea{В \textbf{НОД} каждого простого множителя берём \textbf{минимальную} степень; в \textbf{НОК} --- \textbf{максимальную}. Правило действует для любого конечного числа положительных целых чисел.}

\textbf{Пример из занятия.} Разложим $72$ и $120$:
\[
\begin{aligned}
72&=2^3\cdot3^2\cdot5^0,\\[-1pt]
120&=2^3\cdot3^1\cdot5^1.
\end{aligned}
\]
\[
\begin{aligned}
\gcdop(72,120)&=2^3\cdot3^1\cdot5^0=\mathbf{24},\\
\lcmop(72,120)&=2^3\cdot3^2\cdot5^1=\mathbf{360}.
\end{aligned}
\]
\textit{Проверка:} $24\cdot360=72\cdot120$.

\topic{2. Алгоритм Евклида: сохраняем НОД}
\idea{Если $a,b$ --- положительные целые числа и $a>b$, то вычитание меньшего числа из большего сохраняет общие делители:}
\formula{\gcdop(a,b)=\gcdop(a-b,b).}
\idea{Более общий и особенно полезный приём: для \textbf{любого целого} $k$}
\formula{\gcdop(a,b)=\gcdop(a-kb,b).}
Если $a=qb+r$ и $0\le r<b$, то $\gcdop(a,b)=\gcdop(b,r)$. Повторяем деление с остатком до нулевого остатка; \textbf{последний ненулевой остаток} и есть НОД.

\textbf{Пример.} $\gcdop(252,105)$:
\[
252=2\cdot105+42,\qquad
105=2\cdot42+21,\qquad
42=2\cdot21+0.
\]
Следовательно, $\gcdop(252,105)=\mathbf{21}$.

\subtopic{Числа на фиксированном расстоянии}
\idea{Для любых натуральных $n$ и $k$ верно $\gcdop(n,n+k)=\gcdop(n,k)$.}
\[
\gcdop(n,n+1)=\gcdop(n,1)=\mathbf1.
\]
\textbf{Вывод:} соседние натуральные числа всегда взаимно просты.
\[
\gcdop(n,n+2)=\gcdop(n,2)=
\begin{cases}
2,&\text{если $n$ чётное},\\
1,&\text{если $n$ нечётное}.
\end{cases}
\]

\topic{3. Доказываем взаимную простоту}
Два натуральных числа взаимно просты тогда и только тогда, когда их НОД равен $1$.

\idea{Главный приём: подобрать \emph{кратное} одного выражения так, чтобы при вычитании из другого получилась небольшая константа --- особенно $1$ или $-1$.}

\textbf{Пример 1.} Для любого натурального $n$:
\[
\begin{aligned}
\gcdop(3n+2,3n+5)
 &=\gcdop(3n+2,3)\\
 &=\gcdop(2,3)=\mathbf1.
\end{aligned}
\]
Значит, $3n+2$ и $3n+5$ взаимно просты.

\textbf{Пример 2.} Для любого натурального $n$:
\[
(12n+5)-2(6n+2)=1.
\]
Любой общий делитель исходных выражений делит их разность, равную $1$. Поэтому $\gcdop(12n+5,6n+2)=\mathbf1$.

\subtopic{Важное уточнение о линейных комбинациях}
Если $d$ делит одновременно $a$ и $b$, то $d$ делит также $xa+yb$ для любых целых $x,y$. \textbf{Но произвольная замена одного числа на $xa+yb$ может изменить НОД.}

\textbf{Контрпример:} $\gcdop(12,20)=4$, а $\gcdop(5\cdot12+20,20)=\gcdop(80,20)=20$.

Безусловно сохраняет НОД замена $a\mapsto a+kb$ (или $a\mapsto -a+kb$), где $k$ --- целое. Для доказательства взаимной простоты достаточно найти \textbf{любую} целочисленную комбинацию исходных чисел, равную $\pm1$.

\clearpage
\begin{center}
{\LARGE\bfseries\color{darkaccent}Алгоритм: поиск НОД}\par
\vspace{3pt}
{\small\color{muted}Метод Евклида для положительных целых чисел}\par
\end{center}
\vspace{11mm}
\begin{center}
\begin{tikzpicture}[
    >=Stealth,
    base/.style={draw=darkaccent,thick,align=center,inner sep=7pt,font=\small},
    terminal/.style={base,ellipse,minimum width=46mm,minimum height=11mm,fill=soft!45},
    action/.style={base,rounded corners=2pt,text width=54mm,minimum height=14mm},
    decision/.style={base,diamond,aspect=2.25,text width=27mm,inner sep=3pt},
    arrow/.style={-Stealth,thick,draw=accent},
    every node/.style={font=\small}
]
\node[terminal] (start) at (0,0) {Начало: $a,b>0$};
\node[action] (sort) at (0,-1.9) {Если $a<b$, поменяйте числа местами};
\node[action] (divide) at (0,-4.1) {Разделите $a$ на $b$:\\ $a=qb+r$, $0\le r<b$};
\node[decision] (check) at (0,-7.0) {$r=0$?};
\node[action,text width=42mm] (update) at (-4.0,-10.2) {Замените пару\\ $(a,b)\leftarrow(b,r)$};
\node[action,text width=42mm] (answer) at (4.0,-10.2) {$\gcdop(a,b)=b$};
\node[terminal,minimum width=39mm] (finish) at (4.0,-12.3) {Конец};
\draw[arrow] (start)--(sort);
\draw[arrow] (sort)--(divide);
\draw[arrow] (divide)--(check);
\draw[arrow] (check.west)-| node[pos=0.23,above]{нет} (update.north);
\draw[arrow] (check.east)-| node[pos=0.23,above]{да} (answer.north);
\draw[arrow] (update.west)-- ++(-1.55,0) |- (divide.west);
\draw[arrow] (answer)--(finish);
\end{tikzpicture}
\end{center}
\vfill
\noindent\textbf{Почему алгоритм работает.} После каждого деления НОД пары сохраняется, а меньший положительный аргумент убывает. Поэтому процесс заканчивается, когда остаток станет нулевым.
\clearpage

\topic{4. Выделяем НОД как общий множитель}
\idea{Пусть $d=\gcdop(a,b)$. Тогда $a=dx$ и $b=dy$, где $\gcdop(x,y)=1$. После деления обоих исходных чисел на их НОД остаётся взаимно простая пара.}

\textbf{Пример из занятия:} $16=4\cdot4$ и $36=4\cdot9$; после выделения $d=4$ получаем взаимно простые $4$ и $9$.

\idea{Для любых положительных целых $a,b$ верно тождество}
\formula{\gcdop(a,b)\cdot\lcmop(a,b)=ab.}
Обозначая $m=\lcmop(a,b)$ и пользуясь $a=dx$, $b=dy$, получаем:
\[
d\cdot m=d^2xy\quad\Longrightarrow\quad
\boxed{m=dxy,\qquad xy=\frac md.}
\]

\topic{5. Восстанавливаем пару по НОД и НОК}
\idea{Если известны $d=\gcdop(a,b)$ и $m=\lcmop(a,b)$, сначала вычислите $P=m/d$. Затем найдите \textbf{все} пары натуральных взаимно простых множителей $x,y$, для которых $xy=P$. Восстановите $a=dx$, $b=dy$.}

\textbf{Пример из занятия.} $d=4$, $m=36$:
\[
xy=\frac{36}{4}=9,\qquad (x,y)=(1,9).
\]
Поэтому единственная \textbf{неупорядоченная} пара исходных чисел --- $\pair{4}{36}$.

\textbf{Ещё один пример.} $d=6$, $m=210$. Тогда $xy=35$. Подходят пары $\pair{1}{35}$ и $\pair{5}{7}$, поскольку числа внутри каждой пары взаимно просты. Искомые пары: $\pair{6}{210}$ и $\pair{30}{42}$.

\textbf{Ограничения.} Необходимо, чтобы $d\mid m$; иначе искомых пар нет. Если порядок чисел учитывается, пару с разными числами следует посчитать в двух порядках.

\clearpage

\topic{Тренировочный блок \quad $\cdot$ \quad Путь А}
\textcolor{muted}{10 заданий. Двигайтесь последовательно. Для доказательств записывайте ключевые преобразования; для задач на пары чисел указывайте все варианты.}
\begin{enumerate}[label=\textcolor{accent}{\bfseries\arabic*.},itemsep=7pt]
\item Разложите $48$ и $180$ на простые множители. Найдите их НОД и НОК.
\item По простым множителям найдите $\gcdop(84,126,210)$ и $\lcmop(84,126,210)$.
\item Алгоритмом Евклида вычислите $\gcdop(1105,478)$. Запишите цепочку делений с остатком.
\item Докажите, что для любого натурального $n$ числа $n$ и $n+1$ взаимно просты.
\item Найдите \textbf{все возможные значения} $\gcdop(n,n+6)$ при натуральном $n$. Укажите условие для каждого значения.
\item Найдите $\gcdop(5n+1,5n+6)$ для произвольного натурального $n$ и обоснуйте результат.
\item Докажите, что числа $14n+9$ и $7n+4$ взаимно просты при любом натуральном $n$.
\item Выразите $\gcdop(17n+8,5n+2)$ через НОД небольшого числа и выражения от $n$. Перечислите все возможные значения.
\item Найдите все \textbf{неупорядоченные пары} натуральных чисел $(a,b)$, для которых\[\gcdop(a,b)=6,\qquad\lcmop(a,b)=210.\]
\item Найдите количество и перечислите все \textbf{неупорядоченные пары} натуральных чисел $(a,b)$, для которых\[\gcdop(a,b)=12,\qquad\lcmop(a,b)=360.\]
\end{enumerate}
\vspace{6pt}
\textcolor{muted}{Проверьте: для каждой найденной пары произведение чисел равно произведению её НОД и НОК.}
\clearpage

\topic{Ответы и ориентиры для проверки}
{\small\color{muted}Задания на доказательство содержат главный проверочный шаг.}\par\vspace{7pt}
\renewcommand{\arraystretch}{1.44}
\begin{tabularx}{\linewidth}{@{}>{\bfseries\color{accent}}p{9mm}X@{}}
\toprule
№ & Ответ или ключевой результат \\
\midrule
1 & $\gcdop(48,180)=12$, $\lcmop(48,180)=720$. \\
2 & $\gcdop(84,126,210)=42$, $\lcmop(84,126,210)=1260$. \\
3 & $1$; последние остатки: $25,\;6,\;1,\;0$. \\
4 & $\gcdop(n,n+1)=\gcdop(n,1)=1$. \\
5 & $\gcdop(n,6)\in\{1,2,3,6\}$. Значения по $n\bmod6$:\newline $1$: $n\equiv1,5$; $2$: $n\equiv2,4$; $3$: $n\equiv3$; $6$: $n\equiv0$. \\
6 & $\gcdop(5n+1,5n+6)=\gcdop(5n+1,5)=1$. \\
7 & $(14n+9)-2(7n+4)=1$, поэтому НОД равен $1$. \\
8 & $\gcdop(17n+8,5n+2)=\gcdop(|n-2|,6)$; возможны $1,2,3,6$. \\
9 & Две пары: $\pair{6}{210}$ и $\pair{30}{42}$. \\
10 & Четыре пары: $\pair{12}{360}$, $\pair{24}{180}$, $\pair{36}{120}$, $\pair{60}{72}$. \\
\bottomrule
\end{tabularx}

\vspace{12mm}
{\color{accent}\rule{\linewidth}{0.65pt}}\par
\small\textcolor{muted}{При переборе пар важно одновременно проверить произведение, взаимную простоту множителей $x,y$ и отсутствие повторов при перестановке чисел.}
\end{document}
